\documentclass{article}
 \usepackage{amsmath}
 \usepackage{amsthm}
 \usepackage{amssymb}
 \usepackage{pgfplots}
 \usepackage[utf8]{inputenc}
 \usepackage{amsfonts}
 \usepackage[margin=2.5cm]{geometry}
 \usepackage{graphicx}
 \usepackage[export]{adjustbox}
 \usepackage{fancyhdr}
 \usepackage[portuguese]{babel}
 \usepackage{hyperref}
 \usepackage{lastpage}
 \usepackage{mathtools}
 \setcounter{section}{-1}

 \pagestyle{fancy}
 \fancyhf{}

 \pgfplotsset{compat = 1.18}

 \hypersetup{
     colorlinks,
     citecolor=black,
     filecolor=black,
     linkcolor=black,
     urlcolor=black
 }
 \newtheorem*{def*}{\underline{Defini\c c\~ao}}
 \newtheorem*{theorem*}{\underline{Teorema}}
 \newtheorem*{lemma*}{\underline{Lema}}
 \newtheorem*{prop*}{\underline{Proposi\c c\~ao}}
 \newtheorem{example}{\underline{Exemplo}}
 \newtheorem*{proof*}{\underline{Prova}}
 \renewcommand\qedsymbol{$\blacksquare$}
 \newcommand{\Lin}[1]{Lin_{\mathbb{K}}({#1})}

 \rfoot{P\'agina \thepage \hspace{1pt} de \pageref{LastPage}}

 \begin{document}
 \begin{figure}[ht]
  \minipage{0.76\textwidth}
    \includegraphics[width=4cm]{../icmc.png}
    \hspace{7cm}
    \includegraphics[height=4.9cm,width=4cm]{../brasao_usp_cor.jpg}
  \endminipage  
\end{figure}

\begin{center}
  \vspace{1cm}
  \LARGE
  UNIVERSIDADE DE S\~AO PAULO

  \vspace{1.3cm}
  \LARGE
  INSTITUTO DE CI\^ENCIAS MATEM\'ATICAS E COMPUTACIONAIS - ICMC

  \vspace{1.7cm}
  \Large
  \textbf{Notas de Física II}

  \vspace{1.3cm}
  \large
  \textbf{Renan Wenzel - 11169472}

  \vspace{1.3cm}
  \large
  \textbf{Professor - Luis Gustavo Marcassa}

  \textbf{E-mail: marcassa@ifsc.usp.br}

  \vspace{1.3cm}
  \today
\end{center}

 \newpage
 \textbf{{\Huge Disclaimer}}
 \vspace{5cm}

  {\huge Essas notas não possuem relação com professor algum. 

  Qualquer erro é responsabilidade solene do autor.

Caso julgue necessário, contatar: renan.wenzel.rw@gmail.com}
 \newpage

 \tableofcontents

 \newpage

 \section{Aula 00 - 07/08/2023}
   Avisos sobre o curso (Ler e baixar o pdf no e-disciplinas!!!!); 

\section{Aula 01 - 09/08/2023}
\subsection{Motivações}
\begin{itemize}
  \item Ângulos, velocidade angular e aceleração angular;
  \item Energia em sistemas com rotação.
\end{itemize}
\subsection{Rotação}
  Antes de qualquer coisa, convencionamos o sentido antihorário como aquele em que \(\Delta \theta >0\) e
o sentido horário como o que \(\Delta \theta <0.\) Uma volta completa em torno do círculo é dada pela versão com \(2\pi\) da fórmula do arco de círculo
 \(\Delta S_{i} = r_{i}\Delta \theta = 2\pi r_{i}\) e, com isso, a variação do ângulo em uma volta completa é dada por 
   \[
     \Delta \theta = \frac{S_{i}}{r_{i}} = \frac{2\pi r_{i}}{r_{i}} = 2\pi rad.
   \]
  Um dos assuntos de importância para nós é o estudo da variação temporal do ângulo. Definimos, nessa lógica, a 
velocidade angular média por 
  \[
    \omega _{med} = \frac{\Delta \theta }{\Delta t}.
  \]
  De modo análogo ao que vimos com cinemática, existe também a velocidade angular instantânea, obtida tomando o limite: 
    \[
      \omega = \lim_{\Delta t\to 0}\frac{\Delta \theta }{\Delta t} = \frac{d\theta }{dt}.
    \]
    Observa-se de cara que, se \(\omega >0, \theta \) aumenta e, se \(\omega <0, \theta \) diminui. Assim como antes,
  precisamos ver, também, a unidade. Em cinemática, a unidade de velocidade era metro por segundo. Dessa vez, já que
  o ângulo move-se em radianos, mas há outras unidades, como a revolução e o grau. Logo, as unidades de \(\omega \) podem ser \([\omega ] = \frac{radianos}{tempo} = \frac{rad}{s}, \frac{graus}{s}, \frac{\text{revolução}}{s},\)
  em que \(1\text{revolução} = 2\pi rad = 360\deg\)

    Por exemplo, se um CD roda a 3000rpm, pode-se expressar essa velocidade de rotação como 
      \[
        \omega = 3000rpm = \frac{3000 \cdot 2\pi}{60} = \frac{600}{6}\pi = 100\pi \frac{rad}{s}.
      \]

  Analogamente, é possível analisar a variação da própria velocidade angular com o tempo, resultando na chamada
  acelerações angulares média e instantânea: 
    \[
      \alpha_{med} = \frac{\Delta \omega }{\Delta t}\quad \alpha  = \frac{d\omega }{dt} = \frac{d}{dt}\biggl(\frac{d\theta }{dt}\biggr) = \frac{d^{2}\theta }{dt^{2}}.
    \]
    A unidade dessa grandeza, novamente, similar à versão linear dela, será dada em \([\alpha ]= \frac{radiano}{s^{2}}\), ou \([\alpha ]=\frac{grau}{s^{2}}\), etc. Nessas
  situações todas, se \(\alpha >0, \omega \) aumenta e, se \(\alpha <0, \omega \) diminui. 

    Agora, suponha que \(\alpha \) é constante. Todos os processos de movimento uniformemente acelerado são válidos aqui também:
    \begin{table}[h!]
    \centering
    \begin{tabular}{|c|c|c|}
        \hline
        & \textbf{Variáveis angulares} & \textbf{Variáveis escalares} \\
        \hline
        \textbf{Posição} & $\theta(t) = \theta_{0} + \omega_{0}t + \frac{\alpha t^{2}}{2}$ & $s(t) = R\theta(t)$ \\
        \hline
        \textbf{Velocidade} & $\omega(t) = \omega_{0} + \alpha t = \frac{d\theta }{dt}$ & $v(t) = v_{0} + at = \frac{dx}{dt}$ \\
        \hline
        \textbf{Aceleração} & $\alpha(t) = \frac{d\omega(t)}{dt} = \frac{d^2\theta(t)}{dt^2}$ & $|\vec{a}(t)| = \frac{dv(t)}{dt} = R\alpha(t),\quad |\vec{a}_{cp}| = \frac{v^2}{R}$ \\
        \hline
        \textbf{Torricelli} & $\omega^{2}(t) = \omega_{0}^{2} + 2\alpha \Delta \theta $ & $v^{2} = v_{0}^{2} + 2a\Delta s.$ \\
        \hline
    \end{tabular}
    \caption{Resumo movimento circular.}
    \label{tab:my_label}
  \end{table}

\begin{example}
  Suponha que há um CD que começa no repouso. Ele começa a girar, indo de 0 a 500rpm em 5.5s. Pergunta-se:
 \begin{itemize}
   \item[a)] Quanto vale \(\alpha \)?
     \item[b)] Quantas voltas o CD dá em 5.5s?
       \item[c)] Qual é a distância percorrida por uma ponta a 6cm do eixo de rotação?
 \end{itemize}

 \textbf{Soluções:}
 \begin{itemize}
   \item[a)] Temos \(\omega (0) = 0, \omega (5.5) = 500rpm.\) Segue que 
     \[
       \omega(t) = \omega_{0} + \alpha t \Rightarrow \alpha  = \frac{\omega (t)}{t} = \frac{500 \cdot 2\pi}{5.5 \cdot 60}\approx 9.52 \frac{rad}{s^{2}}
     \]

    \item[b)] Aplicamos o Torricelli angular com os dados que temos: 
      \[
        \omega^{2} = 2\alpha \Delta \theta \Rightarrow \delta \theta = \frac{\omega^{2}}{2\alpha }\approx 144 rad \Rightarrow \frac{144}{2\pi}rad\approx 23\text{rotações}.
      \]

    \item[c)] Por fim, multiplicando a variação do ângulo pelo raio, obtemos 
      \[
        \Delta S_{i} = r\Delta \theta = 6 \cdot 10^{-2}\cdot 144\approx 8.65m.
      \]
 \end{itemize}
\end{example}

  Olhando de forma cautelosa a fórmula de arco de circulo, podemos derivá-la com respeito ao tempo utilizando o que vimos até agora: 
    \[
      \frac{dS_{i}}{dt} = V_{t} = r_{i}\frac{d\theta }{dt} = r_{i}\omega.
    \]
  Essa derivação resulta em uma velocidade linear, que também pode ser derivada a fim de obter uma aceleração linear 
    \[
      \frac{dV_{t}}{dt} = r_{i}\frac{d\omega }{dt} \Rightarrow a_{t} = r_{i}\alpha.
    \]
    Note a relação entre as duas acelerações que obtivemos, \(a_{c} = \frac{V_{t}^{2}}{r_{i}}= \frac{r_{i}^{2}\omega^{2}}{r_{i}} = r_{i}\omega^{2}.\)

\subsection{Energia Cinética de Rotação}
  A energia cinética, como vista previamente, é dada por 
    \[
      \mathcal{K} = \frac{1}{2}m_{i}v_{i}^{2}.
    \]
  Agora, imagine um corpo discreto (formado por vários pontos). Somemos as energias deles, tal que a energia cinética total é 
    \[
      \mathcal{K}_{T} = \sum\limits_{}^{}\frac{1}{2}m_{i}v_{i}^{2}.
    \]
  Mas, sabemos que \(v_{i} = r_{i}\omega, \) tal que 
    \[
      \mathcal{K} = \frac{1}{2}\sum\limits_{}^{}m_{i}r_{i}^{2}\omega^{2} = \frac{1}{2}\biggl[\sum\limits_{}^{}m_{i}r_{i}^{2}\biggr]\omega^{2}
    \]
  Chamemos o termo em colchete de momento de inércia, denotado por \(I:= \sum\limits_{}^{}m_{i}r_{i}^{2} = \sum\limits_{}^{}I_{i}\). Logo, 
    \[
      \boxed{\hypertarget{kin_en}{\mathcal{K}_{T} = \frac{1}{2}I\omega^{2}.}}
    \]
\end{document}
